\scanpage{47}
\margin{46}
\mspage{-3.4-}

space. Then there is a $\mathrm{PGL}(n)$-linear isomorphism:
\[
(*)\qquad U_R\cong\mathrm{PGL}(n)\times\mathbb{A}^{n\nu-n}
\]
Hence $U_R$ is a globally trivial principal fibre bundle with respect to the action of $\mathrm{PGL}(n)$, with base space $\mathbb{A}^{n\nu-n}$.

\uline{Proof}: We shall construct morphisms from the left hand side of $(*)$ to the right, and vice versa, leaving it to the reader to check that they are mutually inverse. The morphisms will be obviously $\mathrm{PGL}(n)$-linear. First of all, we identify $\mathbb{A}^{n\nu-n}$ with the scheme:
\[
\mathbb{A}_R=\operatorname{Spec}
\frac{\mathbb{Z}[x_0^{(n+1)},\ldots,x_n^{(n+1)};\ x_0^{(n+2)},\ldots,x_n^{(n+2)};\ldots;\ x_0^{(n+\nu)},\ldots,x_n^{(n+\nu)}]}
{\text{Ideal generated by }\{x_i^{(n+k)}-1\mid i\in S_k\}}
\]
since both are affine spaces of dimension $n\nu-n$. Then a morphism
\[
\unsure{S_R|s_R}:\mathbb{A}_R\longrightarrow U_R
\]
is defined by the collection of composite morphisms
\diagram{
\begin{tikzcd}[column sep=small]
\unsure{S_i|s_i}: \mathbb{A}_R \arrow[r,"\unsure{S|s}"] & U_R\subset(\mathbb{P}_n)^{n+\nu+1} \arrow[r,"p_i"] & \mathbb{P}_n
\end{tikzcd}
}
and $\unsure{S_i|s_i}$ is defined by the conditions (cf. EGA 2, 4.2.3):
\[
\begin{aligned}
\unsure{S_i|s_i}^*(\underline{o}(1))&\cong\underline{o}_{\mathbb{A}}\\
\unsure{S_i|s_i}^*(X_j)&=\delta_{ij}\quad\text{if }0\leq i\leq n\\
\unsure{S_i|s_i}^*(X_j)&=x_j^{(i)}\quad\text{if }n+1\leq i\leq m
\end{aligned}
\]
